% Generated draft: requires mathematical and source review.
\documentclass[11pt]{article}
\usepackage[margin=25mm]{geometry}
\usepackage{fontspec}
\setmainfont{Latin Modern Roman}
\usepackage{amsmath,amssymb,mathtools,unicode-math}
\setmathfont{Latin Modern Math}
\usepackage{enumitem,needspace,xurl,hyperref}
\setlength{\parindent}{0pt}
\setlength{\parskip}{5pt}
\setlength{\emergencystretch}{3em}
\newcommand{\N}{\mathbb{N}}
\newcommand{\Z}{\mathbb{Z}}
\newcommand{\Q}{\mathbb{Q}}
\newcommand{\R}{\mathbb{R}}
\newcommand{\C}{\mathbb{C}}
\newcommand{\F}{\mathbb{F}}
\newcommand{\A}{\mathbb{A}}
\newcommand{\PP}{\mathbb{P}}
\newcommand{\E}{\mathbb{E}}
\newcommand{\eps}{\varepsilon}
\newcommand{\dd}{\,\mathrm{d}}
\begin{document}
\section*{UnsolvedMath ID: 2307014}
\textbf{Source problem code:} catalogue entry
\subsection{Definitions and mathematical statement}
Let $\gamma$ be an oriented $C^1$ plane curve and let $f:\gamma\to\C$ be continuous. Wherever the integral is well defined, set
\[
 F(z)=\int_\gamma\frac{f(t)}{t-z}\,dt\qquad(z\notin\gamma).
\]
For example, compact rectifiable curves and continuous densities give the ordinary line integral off the curve. The source does not explicitly state compactness, regularity of parametrization, or convergence assumptions for a noncompact curve.

\textbf{Boundary-limit question (A. Baernstein).} Does $F$ have finite non-tangential boundary values at almost every point of $\gamma$, with almost everywhere understood with respect to arclength? A usual non-tangential approach from a complementary domain $\Omega$ requires $z\to\xi\in\partial\Omega$ with $|z-\xi|\leq A\operatorname{dist}(z,\partial\Omega)$ for a fixed $A>1$. For a regular Jordan curve this gives the respective interior and exterior questions; the source does not specify sides or how to treat self-intersections in its first part.

\textbf{Function-class question.} Suppose $\gamma$ is a Jordan curve, let $\Omega$ be its bounded interior, and let $\phi:\mathbb D\to\Omega$ be conformal. For $h=F\circ\phi$, does there exist $0<p<1$ such that
\[
 \sup_{0<r<1}\frac1{2\pi}\int_{-\pi}^{\pi}|h(re^{i\theta})|^p\,d\theta<\infty,
\]
that is, $h\in H^p$? If this Hardy-class conclusion is unavailable, does at least $h\in N$, the analytic class of bounded Nevanlinna characteristic, hold, meaning
\[
 \sup_{0<r<1}\frac1{2\pi}\int_{-\pi}^{\pi}\log^+|h(re^{i\theta})|\,d\theta<\infty?
\]
The natural reading allows $p$ to depend on the curve and density, but the source does not specify whether a universal exponent is sought. It is also unclear whether the original $C^1$ hypothesis is retained in this Jordan-curve part; for an arbitrary nonrectifiable Jordan curve the meaning of the original line integral itself requires an extra convention.

The source describes a long-studied old problem, reports affirmative answers under slightly stronger smoothness assumptions on either $\gamma$ or $f$, and highlights Havin's contribution without stating its hypotheses. The supplied summary says the question was open in the 2018 edition and no definitive resolution was located. These are historical/supplied status claims, not an independent present-day literature assessment.
\subsection{Short English statement}
For a continuous density on a $C^1$ curve, ask whether its Cauchy integral has non-tangential limits almost everywhere. On a Jordan-curve interior, ask whether conformal pullback lies in some $H^p$, $0<p<1$, or at least has bounded Nevanlinna characteristic. Curve scope and integral conventions need clarification.
\subsection{Sources}
\begin{itemize}
\item Local UnsolvedMath registry, permanent ID 2307014, source code catalogue entry, ``Research Problems in Function Theory --- Problem 7.14''; attribution: A. Baernstein. Supplied status open, classification OPEN-TRIAGE, updated 2026-08-06.
\item Hayman, \emph{Research Problems in Function Theory} (2018), Problem 7.14, \url{https://arxiv.org/abs/1809.07200}. The record reports source-TeX extraction and access on 2026-07-29; the 2018 update was retained for status triage. Current status and rights require review; the record describes publicly available arXiv source TeX.
\item The statement refers to extensive older work, especially Russian work, and singles out Havin for an outstanding contribution under stronger smoothness assumptions. No complete citation or precise theorem is supplied.
\end{itemize}
\end{document}
